% TOP_PROBLEM: {"id": 30007023, "title": "Kähler log abundance assuming logarithmic Iitaka", "category_id": 5, "category": {"id": 5, "name": "algebraic_geometry", "display_name": "Algebraic Geometry", "description": "Geometric objects defined by polynomial equations.", "slug": "algebraic-geometry", "order_index": 5, "created_at": "2026-07-31T15:26:25.671Z"}, "status": "open"}
% FIRST_POSTED: 2026-10-06
% ATTEMPT_STATUS: solved
% OPENAI_CLAIM: {"schema_version":1,"source_repo":"https://github.com/openai/math","source_commit":"adc7f1241b42e322a6451854ab7e4b4c146bf78a","release_date":"2026-10-06","match":"full","resolution":"proved","independently_reviewed":false,"second_pass":"agreed","families":[{"family":"034","title":"Log abundance for compact Kähler spaces under logarithmic Iitaka subadditivity","subject":"Algebraic and complex geometry","manuscripts":[{"title":"Log abundance for compact Kähler spaces under logarithmic Iitaka subadditivity","pdf_path":"preprints/Log-abundance-for-compact-Kahler-spaces-under-logarithmic-Iitaka-subadditivity-October-4-2026/main.pdf","readme_path":"preprints/Log-abundance-for-compact-Kahler-spaces-under-logarithmic-Iitaka-subadditivity-October-4-2026/README.md","date":"2026-10-04","theorem_ref":"thm:main (sections/01-introduction.tex:56)","statement_sources":[{"path":"preprints/Log-abundance-for-compact-Kahler-spaces-under-logarithmic-Iitaka-subadditivity-October-4-2026/build/sections/01-introduction.tex","line":56,"label":"thm:main"}]}],"lean":null,"reasoning_trace":null}]}
\documentclass[11pt]{article}
\usepackage[margin=25mm]{geometry}
\usepackage{fontspec}
\setmainfont{Latin Modern Roman}
\usepackage{amsmath,amssymb,mathtools}
\usepackage{unicode-math}
\setmathfont{Latin Modern Math}
\usepackage{enumitem}
\usepackage{xurl,hyperref,bookmark}
\hypersetup{colorlinks=true,urlcolor=blue}
\setcounter{secnumdepth}{0}
\setlength{\parindent}{0pt}
\setlength{\parskip}{5pt}
\setlength{\emergencystretch}{3em}
\begin{document}
\section*{Kähler log abundance assuming logarithmic Iitaka}

\subsection{OpenAI's claim}
Assuming the stated logarithmic Iitaka inequality, OpenAI claims semiampleness of nef rational log adjoints on compact Kähler log-canonical spaces in arbitrary dimension. The cited manuscript argues for that result; its relation to this problem is specified below.

\subsection{Mathematical statement of the claimed result}
\paragraph{Mathematical claim 1.} Assume logarithmic Iitaka subadditivity for every connected-fiber surjective morphism $f:U\to V$ of smooth connected projective complex varieties with reduced simple-normal-crossing divisors $D_U,D_V$ satisfying $\operatorname{Supp}(f^*D_V)\subseteq\operatorname{Supp}D_U$: for a very general smooth fiber $F$ and $D_F=D_U|_F$,
\[
\kappa(U,K_U+D_U)\ge\kappa(F,K_F+D_F)+\kappa(V,K_V+D_V).
\]
Under that premise, OpenAI claims that every normal irreducible compact Kähler complex space $X$ with effective rational $\Delta$, log canonical pair $(X,\Delta)$ and analytically nef rational Cartier $D=K_X+\Delta$ has a positive Cartier multiple generated everywhere by its global sections:
\[
H^0(X,\mathcal O_X(mD))\otimes\mathcal O_X\twoheadrightarrow\mathcal O_X(mD).
\]
The premise covers all dimensions, uses $\kappa=-\infty$ for empty positive systems, and is a hypothesis of this Kähler result.

\subsection{Outline of the argument}
Manuscript 1 (family 034, 2026-10-04): The induction retains both the nef line and the effective negative part through a restricted birational program. Boundary generation, signed-support root constructions, Hodge-module lifting, and compact cycle deformations produce nonvanishing; affine changes in negative multiplicities make the selected wall crossings finite.

\subsection{Scope relative to this problem}
Manuscript 1 (family 034, 2026-10-04): The problem selects the mathematical statement summarized here, including its stated hypotheses and parameter range. The cited theorem gives a claimed proof of that assertion in full.

\subsection{Formal verification}
Family 034 has no Lean scope page in this snapshot; no formalization of these manuscript results is listed.

This site has not executed the Lean files or checked their dependencies. These descriptions report the scope of the linked declarations, rather than a new verification of the proofs.

\subsection{Status caveat}
OpenAI attributes these results to an internal model. Its release notice cautions that results lacking full formalization can contain mistakes. This site has not independently assessed the proofs; a solved label records the scope of the reported claim.

\subsection{Sources}
\begin{itemize}
\item OpenAI, \emph{Log abundance for compact Kähler spaces under logarithmic Iitaka subadditivity}, 2026-10-04; citation key \texttt{OAI:Log-abundance-for-compact-Kahler-spaces-under-logarithmic-Iitaka-subadditivity-October-4-2026}. \href{https://github.com/openai/math/blob/adc7f1241b42e322a6451854ab7e4b4c146bf78a/preprints/Log-abundance-for-compact-Kahler-spaces-under-logarithmic-Iitaka-subadditivity-October-4-2026/main.pdf}{PDF (GitHub)}; \href{https://github.com/openai/math/raw/adc7f1241b42e322a6451854ab7e4b4c146bf78a/preprints/Log-abundance-for-compact-Kahler-spaces-under-logarithmic-Iitaka-subadditivity-October-4-2026/main.pdf}{PDF (direct)}; \href{https://github.com/openai/math/blob/adc7f1241b42e322a6451854ab7e4b4c146bf78a/preprints/Log-abundance-for-compact-Kahler-spaces-under-logarithmic-Iitaka-subadditivity-October-4-2026/README.md}{Manuscript page and citation}; \href{https://github.com/openai/math/blob/main/preprints/Log-abundance-for-compact-Kahler-spaces-under-logarithmic-Iitaka-subadditivity-October-4-2026/main.pdf}{Latest version}; Statement source: \href{https://github.com/openai/math/blob/adc7f1241b42e322a6451854ab7e4b4c146bf78a/preprints/Log-abundance-for-compact-Kahler-spaces-under-logarithmic-Iitaka-subadditivity-October-4-2026/build/sections/01-introduction.tex\#L56}{thm:main}.
\item \href{https://github.com/openai/math/blob/adc7f1241b42e322a6451854ab7e4b4c146bf78a/CONTENTS.md}{OpenAI catalogue entry}.
\end{itemize}
\end{document}
